\addcontentsline{toc}{section}{Введение}
\section*{Введение}

В ходе курсовой работы требуется сформировать калькулятор «большой» конечной арифметики
\(\langle Z^8_8;\ +,\ * \rangle\) для четырех действий
\((+,\ -,\ *,\ \div)\) на основе «малой» конечной арифметики, где задано
правило «+1» и выполняются свойства коммутативности \((+,\ *)\),
ассоциативности \((+,\ *)\), дистрибутивности \(*\) относительно \(+\),
заданы аддитивная единица «a» и мультипликативная единица «b», а также
выполняется свойство: для \((\forall x)\ x * a = a.\)


Правило "$+1$" определяет переход от текущего символа к следующему и задано в соответствии с таблицей ~\ref{tab:plus_one}.

\begin{table}[H]
    \centering
    \caption{Правило "$+1$"}
    \label{tab:plus_one}
    \begin{tabular}{|c|c|c|c|c|c|c|c|c|}
        \hline
        $x$ & a & b & c & d & e & f & g & h \\
        \hline
        $x + 1$ & b & c & e & g & d & h & f & a \\
        \hline
    \end{tabular}
\end{table} 
\newpage